\section*{Chapitre \ref{ch:g7:fractions} --- Fractions~: comparer et additionner}

\begin{solution}{exo:g7:fractions:1}
$\dfrac{2}{5} = \dfrac{8}{20}$~; \quad
$\dfrac{18}{24} = \dfrac{3}{4}$~; \quad
$\dfrac{7}{3} = \dfrac{28}{12}$~; \quad
$\dfrac{5}{9} = \dfrac{20}{36}$.
\end{solution}

\begin{solution}{exo:g7:fractions:2}
$\dfrac{12}{16} = \dfrac34$~; \quad
$\dfrac{30}{45} = \dfrac23$~; \quad
$\dfrac{27}{36} = \dfrac34$~; \quad
$\dfrac{48}{12} = 4$.
\end{solution}

\begin{solution}{exo:g7:fractions:3}
$\dfrac{15}{35} = \dfrac37$ et $\dfrac{21}{49} = \dfrac37$~: égales.

$\dfrac{16}{28} = \dfrac47$ et $\dfrac{20}{36} = \dfrac59$~: pour
comparer, $\frac47 = \frac{36}{63}$ et $\frac59 = \frac{35}{63}$ ---
pas égales.
\end{solution}

\begin{solution}{exo:g7:fractions:4}
$\dfrac{7}{11} < \dfrac{9}{11}$ (même \omterm{def:g4:fractions:def}{dénominateur}).

$\dfrac56 = \dfrac{15}{18} > \dfrac{11}{18}$.

$\dfrac{13}{12} > 1$ tandis que $\dfrac{19}{20} < 1$, donc
$\dfrac{13}{12} > \dfrac{19}{20}$.
\end{solution}

\begin{solution}{exo:g7:fractions:5}
$\dfrac59 + \dfrac29 = \dfrac79$~; \quad
$\dfrac{11}{7} - \dfrac47 = \dfrac77 = 1$~; \quad
$\dfrac58 + \dfrac78 = \dfrac{12}{8} = \dfrac32$.
\end{solution}

\begin{solution}{exo:g7:fractions:6}
$\dfrac23 + \dfrac56 = \dfrac46 + \dfrac56 = \dfrac96 = \dfrac32$.

$\dfrac{7}{10} - \dfrac25 = \dfrac{7}{10} - \dfrac{4}{10} =
\dfrac{3}{10}$.

$\dfrac34 + \dfrac{5}{12} = \dfrac{9}{12} + \dfrac{5}{12} =
\dfrac{14}{12} = \dfrac76$.
\end{solution}

\begin{solution}{exo:g7:fractions:7}
$3 - \dfrac54 = \dfrac{12}{4} - \dfrac54 = \dfrac74$~; \quad
$1 + \dfrac38 = \dfrac{11}{8}$~; \quad
$2 - \dfrac76 = \dfrac{12}{6} - \dfrac76 = \dfrac56$.
\end{solution}

\begin{solution}{exo:g7:fractions:8}
$5 \times \dfrac{3}{10} = \dfrac{15}{10} = \dfrac32$~; \quad
$\dfrac78 \times 4 = \dfrac{28}{8} = \dfrac72$~; \quad
$\dfrac23$ de $45$ vaut $\dfrac{2 \times 45}{3} = \dfrac{90}{3} = 30$.
\end{solution}

\begin{solution}{exo:g7:fractions:9}
Ensemble~: $\dfrac14 + \dfrac38 = \dfrac28 + \dfrac38 = \dfrac58$ de
la pizza. \omterm{def:g3:division:remainder}{Reste}~: $1 - \dfrac58 = \dfrac38$.
\end{solution}

\begin{solution}{exo:g7:fractions:10}
Versé~: $2 \times \dfrac16 = \dfrac26 = \dfrac13$ L. Restant~:
\[
\frac34 - \frac13 = \frac{9}{12} - \frac{4}{12} = \frac{5}{12}
\text{ L}.
\]
\end{solution}

\begin{solution}{exo:g7:fractions:11}
Ajouter la quantité positive $\frac12$ à $\frac35$ doit donner un
résultat \emph{plus grand} que $\frac35$. Or $\frac47 < \frac35$
(en effet $\frac47 = \frac{20}{35}$ et $\frac35 = \frac{21}{35}$)~: la
\omterm{def:g1:addition:def}{somme} annoncée est plus petite que l'un de ses termes --- impossible.
\end{solution}

\begin{solution}{exo:g7:fractions:12}
Étape par étape~:
\[
\frac12 + \frac14 = \frac34, \qquad
\frac34 + \frac18 = \frac68 + \frac18 = \frac78, \qquad
\frac78 + \frac1{16} = \frac{14}{16} + \frac1{16} = \frac{15}{16}.
\]
Les \omterm{def:g1:addition:def}{sommes} partielles $\frac12, \frac34, \frac78, \frac{15}{16},
\dots$ se rapprochent à chaque fois de \omterm{def:g3:division:half}{moitié} vers $1$~: à chaque
étape, la part manquante est coupée en deux. Les \omterm{def:g1:addition:def}{sommes} s'approchent
de $1$ sans jamais l'atteindre.
\end{solution}

\begin{solution}{pb:g7:fractions:1}
\textbf{1.} (a) $\frac12 + \frac14 + \frac14 = \frac24 + \frac14
+ \frac14 = \frac44 = 1$~: correcte. (b) $\frac14 + \frac18 +
\frac18 + \frac12 = \frac28 + \frac18 + \frac18 + \frac48 =
\frac88 = 1$~: correcte.

\textbf{2.} $\frac12 + \frac14 + \frac18 = \frac48 + \frac28 +
\frac18 = \frac78$~: il manque une croche.

\textbf{3.} Une mesure vaut $1 = \frac{16}{16}$~: seize doubles
croches. Une croche vaut $\frac18 = \frac{2}{16}$~: deux doubles
croches.

\textbf{4.} Blanche pointée~: $\frac12 + \frac14 = \frac34$. Noire
pointée~: $\frac14 + \frac18 = \frac28 + \frac18 = \frac38$.

\textbf{5.} La blanche pointée vaut $\frac34$ (question 4)~:
exactement une mesure de valse. Autres remplissages, par exemple~:
noire $+$ noire $+$ noire ($\frac14 + \frac14 + \frac14 =
\frac34$), ou noire $+$ croche $+$ croche $+$ noire
($\frac14 + \frac18 + \frac18 + \frac14 = \frac{2 + 1 + 1 +
2}{8} = \frac68 = \frac34$).

\textbf{6.} $\frac38 + \frac18 + \frac14 = \frac38 + \frac18 +
\frac28 = \frac68 = \frac34$~: il manque une noire ($\frac14$).

\textbf{7.} Croche pointée~: $\frac18 + \frac{1}{16} =
\frac{2}{16} + \frac{1}{16} = \frac{3}{16}$. Noire~:
$\frac{4}{16}$. La noire est la plus longue.

\textbf{8.} $\frac12 + \frac18 = \frac48 + \frac18 = \frac58$,
c'est-à-dire $\frac{10}{16}$~: dix doubles croches.

\textbf{9.} $1 - \frac{11}{16} = \frac{16}{16} - \frac{11}{16}
= \frac{5}{16}$ de la mesure est du silence.

\textbf{10.} Les deux exemples totalisent
$\frac{2 + 1 + 1}{16} = \frac{4}{16}$~: correct. Tous les rythmes
remplissant quatre doubles croches avec des morceaux de $2$ et de
$1$~:
\[
2{+}2, \quad 2{+}1{+}1, \quad 1{+}2{+}1, \quad 1{+}1{+}2, \quad
1{+}1{+}1{+}1 :
\]
cinq rythmes.

\textbf{11.} Un rythme remplissant $n$ doubles croches se termine
soit par une double croche --- et ce qui précède remplit $n - 1$ ---
soit par une croche --- et ce qui précède remplit $n - 2$. Chaque
rythme est ainsi compté exactement une fois, donc
$\text{rythmes}(n) = \text{rythmes}(n-1) + \text{rythmes}(n-2)$.
Tableau~:
\begin{center}
\renewcommand{\arraystretch}{1.2}
\begin{tabular}{c|cccccccc}
$n$ & $1$ & $2$ & $3$ & $4$ & $5$ & $6$ & $7$ & $8$ \\
\hline
rythmes & $1$ & $2$ & $3$ & $5$ & $8$ & $13$ & $21$ & $34$ \\
\end{tabular}
\end{center}
Une blanche ($8$ doubles croches) se bat de $34$ façons. ($n = 4$
redonne les cinq rythmes de la question 10.)

\textbf{12.} Chaque nombre est la \omterm{def:g1:addition:def}{somme} des deux précédents~:
$3 + 5 = 8$, $5 + 8 = 13$, $8 + 13 = 21$, $13 + 21 = 34$ --- la
règle démontrée à la question 11, vieille de mille ans et connue
aujourd'hui sous le nom de règle de Fibonacci.

\textbf{13.} $\frac16 + \frac{1}{12} = \frac{2}{12} +
\frac{1}{12} = \frac{3}{12} = \frac14$~: la même durée que deux
croches ordinaires ($\frac18 + \frac18 = \frac14$). L'échange du
swing est correct.

\textbf{14.} Chaque note du triolet dure $\frac{1}{12}$~:
vérification, $3 \times \frac{1}{12} = \frac{3}{12} = \frac14$
(\cref{prop:g7:fractions:mult}). Pour une blanche~: chacune dure
$\frac16$, puisque $3 \times \frac16 = \frac36 = \frac12$.

\textbf{15.} Noire pointée $+$ noire $+$ noire~: en croches,
$3 + 2 + 2 = 7$ croches $= \frac78$~: correct. Les blanches et les
noires valent $4$ et $2$ croches --- deux nombres \emph{\omterm{def:g3:numbers:evenodd}{pairs}}.
Toute \omterm{def:g1:addition:def}{somme} de \omterm{def:g3:numbers:evenodd}{nombres pairs} est paire, alors qu'une mesure à $7/8$
demande $7$ croches, un \omterm{def:g3:numbers:evenodd}{nombre impair}~: c'est impossible. Pour
remplir une mesure impaire, il faut faire apparaître quelque chose
d'\omterm{def:g3:numbers:evenodd}{impair} --- une croche ou une note pointée.
\end{solution}
